\section{Open questions (5)}

\subsection*{Q1. Hardware-noise sensitivity of the Eq.~12 iteration}
\textbf{Basis.} The replication ran only in ideal Aer statevector; the paper's
Rabi oscillation peaks at $k = \pi/(2 c_n)$, which for small $c_n$ means many gate
applications, and each $U_G U_O$ is a deep multi-controlled-phase circuit on $n$
qubits. Neither the paper nor this replication quantified how much noise erodes
the 0.9991 peak. The naive-vs-constructed 0.30 vs 0.9991 gap could shrink
dramatically under realistic error rates.\\
\textbf{Next steps.} Rerun \texttt{code/generalized\_grover\_v2.py} with
\texttt{AerSimulator(method='density\_matrix')} and a \texttt{NoiseModel} built
from a fake IBM backend (FakeManila, FakeGuadalupe) or hand-set depolarizing
$p \in \{10^{-4}, 10^{-3}, 10^{-2}\}$; sweep peak $P_T$ vs.\ $(p, N, M, D)$; plot
the noise threshold at which the constructed-vs-naive advantage disappears. Free
endpoint: local Qiskit Aer only.

\subsection*{Q2. Extension to spatial search / quantum walks}
\textbf{Basis.} The paper's $H = P_S + P_T$ structure is close in spirit to a
continuous-time quantum walk with a target-marker Hamiltonian, but $P_S$ here is
abstract (random orthonormal vectors), not tied to a graph Laplacian. Whether the
Eq.~12 trick recovers optimal spatial-search scaling on e.g.\ hypercube, complete
graph, or 2D lattice is untested and would connect this work to the
Childs--Goldstone (2004) and Chakraborty--Novo--Ambainis--Omar (2016) lines.\\
\textbf{Next steps.} Replace $P_S$ by the projector onto the low-energy subspace
of a graph Laplacian $L$ (hypercube, then 2D torus); recompute
$c_n = \mathrm{svd}(P_T P_S)$, Eq.~12 initial state, and $P_T(t)$ curve; compare
optimal query count vs.\ $\sqrt{D}$ and published spatial-search bounds. Local
NetworkX + Qiskit; no paid endpoints.

\subsection*{Q3. ML-optimized phase schedules vs.\ analytic Eq.~12 in the near-degenerate regime}
\textbf{Basis.} The paper picks the analytic Eq.~12 preparation and standard $-1$
phase inversions. But when several $c_n$ are close (as in our $D=32$ spectrum:
$0.620, 0.407, 0.193, 0.032$ --- the last two are small and would take many
iterations), a QAOA-style variational schedule with per-iteration phases might
amplify the target subspace faster or with better noise tolerance. The
replication did not test any variational alternative.\\
\textbf{Next steps.} Formulate a $p$-layer variational ansatz
$U(\gamma_p, \phi_p)\cdots U(\gamma_1, \phi_1)|\mathrm{init}\rangle$ minimizing
$1 - P_T$; optimize with COBYLA/SPSA using the same $H$ and $P_T$; compare
depth-to-peak vs.\ the analytic Eq.~12 recipe at fixed $n=5$ and $n=7$. All
local Qiskit.

\subsection*{Q4. Composition with QSVT / block-encoding amplitude amplification}
\textbf{Basis.} Modern quantum search unifies Grover, fixed-point, and oblivious
amplitude amplification under Quantum Singular Value Transformation
(Gily\'en--Su--Low--Wiebe 2019). The Byrnes--Forster--Tessler generalization to
multi-source/multi-target may or may not be a special case of QSVT applied to the
block-encoding of $P_T P_S$. If it is, its query complexity is bounded by QSVT
lower bounds; if it isn't, it may have a distinct optimality story. Not addressed
by paper or replication.\\
\textbf{Next steps.} Write the $U_G U_O$ iteration as a QSVT polynomial in the
block-encoding of $P_T P_S$ (whose singular values are the $c_n$); check whether
the Eq.~12 preparation corresponds to a specific QSVT phase-factor sequence via
\texttt{pyqsp} (Martyn--Rossi--Tan--Chuang); compare achievable $P_T(k)$ curves.

\subsection*{Q5. Analog / adiabatic / counter-diabatic protocol on analog simulators}
\textbf{Basis.} The replication's continuous-time simulation $e^{-iHt}$ already
showed perfect Rabi oscillations at $t = \pi/(2 c_1)$ --- an analog Hamiltonian
evolution under $H = P_S + P_T$ \emph{is} a valid analog protocol, but constructing
$H$ physically requires simultaneous multi-body projectors, which is hard. An
adiabatic sweep from a trivially-preparable Hamiltonian to $H$, or a
counter-diabatic driving scheme, might sidestep the Eq.~12 state-preparation
problem entirely and enable analog implementations (Rydberg arrays, trapped ions,
cQED).\\
\textbf{Next steps.} Simulate an adiabatic schedule $H(s) = (1-s) H_0 + s (P_S + P_T)$
with $H_0 = |0\rangle\langle 0|$; use SciPy ODE integration; measure final $P_T$
vs.\ sweep time $T$ and vs.\ $1/\mathrm{gap}^2$; compare to instantaneous Eq.~12
preparation cost. Also check Sels--Polkovnikov counter-diabatic terms. Local SciPy
only.
